% TODO make hw stylesheet
\documentclass[10pt]{article}
\usepackage[letterpaper,top=1in,left=1in,right=1in]{geometry}
\usepackage[dvipsnames,svgnames]{xcolor}
\usepackage[shortlabels]{enumitem}
\usepackage[framemethod=TikZ]{mdframed}
\usepackage{amsmath,amssymb,amsthm}
\usepackage[colorlinks]{hyperref}
\usepackage{multicol}
\usepackage{tabularx}
\usepackage{bbm}

\hypersetup{
    colorlinks=true,
    linkcolor=blue,
    filecolor=magenta,
    urlcolor=magenta,
    pdftitle={MAT 344 HW}
}

\usepackage{mathtools}
\usepackage{thmtools}
\usepackage[textsize=footnotesize]{todonotes}
\usepackage{titling}
\usepackage{cancel}

\reversemarginpar
\mdfdefinestyle{mdbluebox}{roundcorner=10pt,innerbottommargin=9pt,
linecolor=blue,backgroundcolor=TealBlue!5,}
\declaretheoremstyle[headfont=\sffamily\bfseries\color{MidnightBlue},
mdframed={style=mdbluebox},]{thmbluebox}

\mdfdefinestyle{mdbloobox}{frametitlefont=\bfseries,innerbottommargin=8pt,
nobreak=true,backgroundcolor=TealBlue!5,linecolor=blue,}
\declaretheoremstyle[headfont=\bfseries\color{MidnightBlue},
mdframed={style=mdbloobox},headpunct={\\[3pt]},postheadspace=0pt,]{thmbloobox}

\mdfdefinestyle{mdredbox}{frametitlefont=\bfseries,innerbottommargin=8pt,
nobreak=true,backgroundcolor=Salmon!5,linecolor=RawSienna,}
\declaretheoremstyle[headfont=\bfseries\color{RawSienna},
mdframed={style=mdredbox},headpunct={\\[3pt]},postheadspace=0pt,]{thmredbox}

\mdfdefinestyle{mdgreenbox}{linecolor=ForestGreen,backgroundcolor=ForestGreen!5,
linewidth=2pt,rightline=false,leftline=true,topline=false,bottomline=false,}
\declaretheoremstyle[headfont=\bfseries\sffamily\color{ForestGreen!70!black},
mdframed={style=mdgreenbox},headpunct={ --- },]{thmgreenbox}

\mdfdefinestyle{mdorangebox}{linecolor=RedOrange,backgroundcolor=RedOrange!5,
linewidth=2pt,rightline=false,leftline=true,topline=false,bottomline=false,}
\declaretheoremstyle[headfont=\bfseries\sffamily\color{RedOrange!70!black},
mdframed={style=mdorangebox},headpunct={ --- },]{thmorangebox}

\mdfdefinestyle{mdblackbox}{linecolor=black,backgroundcolor=RedViolet!5!gray!5,
linewidth=3pt,rightline=false,leftline=true,topline=false,bottomline=false,}
\declaretheoremstyle[mdframed={style=mdblackbox}]{thmblackbox}

\declaretheoremstyle[spaceabove=6pt,spacebelow=6pt]{boldhead}
\declaretheoremstyle[spaceabove=6pt,spacebelow=6pt,headfont=\normalfont\itshape]{italichead}

\declaretheorem[style=boldhead,name=Problem]{problem}
\declaretheorem[style=boldhead,name=Problem,numbered=no]{problem*}
\declaretheorem[style=boldhead,name=Setup]{setup} % for config geo
\declaretheorem[style=boldhead,name=Example,sibling=problem]{example}
\declaretheorem[style=boldhead,name=Theorem,sibling=problem]{theorem}
\declaretheorem[style=boldhead,name=Corollary,sibling=theorem]{corollary}
\declaretheorem[style=boldhead,name=Proposition,sibling=theorem]{prop}
\declaretheorem[style=boldhead,name=Lemma,numberwithin=problem]{lemma}
\declaretheorem[style=boldhead,name=Theorem,numbered=no]{theorem*}
\declaretheorem[style=boldhead,name=Lemma,numbered=no]{lemma*}
\declaretheorem[style=boldhead,name=Claim,numberwithin=problem]{claim}
\declaretheorem[style=boldhead,name=Claim,numbered=no]{claim*}
\declaretheorem[style=boldhead,name=Remark,numberwithin=problem]{remark}
\declaretheorem[style=boldhead,name=Remark,numbered=no]{remark*}
\declaretheorem[style=boldhead,name=Definition,sibling=theorem]{definition}
\declaretheorem[style=boldhead,name=Definition,numbered=no]{definition*}

\providecommand{\ol}{\overline}
\providecommand{\ceil}[1]{\left\lceil #1 \right\rceil}
\providecommand{\floor}[1]{\left\lfloor #1 \right\rfloor}
\newcommand{\dd}{\mathrm{d}}
\providecommand{\ul}{\underline}
\providecommand{\eps}{\varepsilon}
\providecommand{\half}{\frac{1}{2}}
\providecommand{\CC}{\mathbb C}
\providecommand{\FF}{\mathbb F}
\providecommand{\II}{\mathbb I}
\providecommand{\NN}{\mathbb N}
\providecommand{\QQ}{\mathbb Q}
\providecommand{\RR}{\mathbb R}
\providecommand{\ZZ}{\mathbb Z}
\providecommand{\ind}{\mathbbm 1}
\providecommand{\dg}{^\circ}
\providecommand{\ii}{\item}
\providecommand{\setdiff}{\smallsetminus}
\providecommand{\alert}[1]{{\sffamily\textbf{\textcolor{blue}{#1}}}}
\providecommand{\inv}{^{-1}}
\providecommand{\mb}[1]{\mathbf{#1}}
\newcommand{\what}{\widehat}

\newenvironment{soln}{\begin{proof}[Solution]}{\end{proof}}

\providecommand{\pow}{\operatorname{Pow}}
\providecommand{\vocab}[1]{{\textbf{\textcolor{ForestGreen}{#1}}}}
\providecommand{\lcm}{\operatorname{lcm}}
\providecommand{\abs}[1]{\left\lvert #1\right\rvert}
\providecommand{\smabs}[1]{\lvert #1\rvert}
\providecommand{\innerprod}[1]{\langle\!\langle #1\rangle\!\rangle}
\providecommand{\norm}[1]{\left\| #1\right\|}
\providecommand{\smnorm}[1]{\| #1\|}
\providecommand{\gr}{\operatorname{gr}}
\providecommand{\im}{\operatorname{im}}

\usepackage{mathrsfs}

% place title at top right
\setlength{\droptitle}{-8ex}
\pretitle{\begin{flushright}\Large\bfseries}
\posttitle{\par\end{flushright}}
\preauthor{\begin{flushright}}
\postauthor{\end{flushright}}
\predate{\begin{flushright}}
\postdate{\end{flushright}}

\title{MAT 344 HW06}
\author{\large Aditya Pahuja}
\date{\large Due 03/10}
\begin{document}
\maketitle

\begin{problem}
    [p.161\#1ce]
    Find the poles and residues of the following functions:
    (c) $\frac 1{\sin z}$; (e) $\frac 1{\sin^2z}$.
\end{problem}

\begin{soln}
    The zeroes of $\sin z$ (hence also of $\sin^2 z$)
    are the points at which
    \begin{equation*}
	\sin z = \frac{e^{iz} - e^{-iz}}{2i} = 0
	\iff e^{2iz} = 1
	\iff z = n\pi\text{ for }n\in\ZZ.
    \end{equation*}
    At $z = n\pi$, $\sin z$ locally has the power series representation
    \begin{equation*}
	\sum_{k=1}^\infty \frac{(z-n\pi)^{2k-1}}{(2k-1)!},
    \end{equation*}
    so
    \begin{equation*}
	1 = \lim_{z\to n\pi} \frac{z-n\pi}{\sin(z)} =
	\lim_{z\to n\pi} \frac{(z-n\pi)^2}{\sin^2(z)}
    \end{equation*}
    (with the second equality coming from
    the continuity of squaring).
    This means that for every integer $n$, $n\pi$ is an order $1$ pole of
    $\frac 1{\sin z}$ and an order $2$ pole of $\frac 1{\sin^2z}$.
    The formula for residues that we calculated in class
    implies that the answer to (c) is
    \begin{equation*}
	\operatorname{Res}_{z=n\pi}\left(\frac 1{\sin z}\right)
	= \lim_{z\to n\pi}\frac{z-n\pi}{\sin z} = \boxed 1.
    \end{equation*}
    For (e),
    \begin{align*}
	\operatorname{Res}_{z=n\pi}\left(\frac 1{\sin^2 z}\right)
	&= \lim_{z\to n\pi}
	\frac{\dd}{\dd z}\frac{(z-n\pi)^2}{\sin^2 z}\\
	&= \lim_{z\to 0} \frac{2z\sin^2 z - 2z^2\sin z\cos z}{\sin^4z}\\
	&= \lim_{z\to 0} \frac{2z\sin z - 2z^2\cos z}{\sin^3z}\\
	&= \lim_{z\to 0} \frac{2\sin z - 2z\cos z}{z^2}.
    \end{align*}
    If we write out the power series for $\sin z$
    and $z\cos z$, the linear terms cancel,
    so that $\sin z - z\cos z$ has a zero of order $3$ at $0$,
    meaning that the limit at the end,
    hence the residues of $\frac 1{\sin^2z}$, are $\boxed 0$.
\end{soln}

\pagebreak
\begin{problem}
    [p.161\#cd]
    Evaluate the following integrals by the method of residues:
    \begin{equation*}
	\text{(c) }\int_{-\infty}^\infty
	\frac{x^2 - x + 2}{x^4 + 10x^2 + 9}\dd x,\qquad
	\text{(d) }\int_0^\infty \frac{x^2}{(x^2 + a^2)^3}\dd x,\; a\in\RR.
    \end{equation*}
\end{problem}

\begin{soln}
    (c) Let $\gamma_R$ be the counterclockwise-oriented
    semicircular arc above the real axis of radius $R$ centered at $0$.
    By the residue theorem, if $R$ is large enough (greater than $3$
    should suffice, since the roots of $x^4 + 10x^2 + 9$
    are $\pm i$, $\pm 3i$), then
    \begin{equation*}
	\int_{-R}^R \frac{x^2 - x + 2}{x^4 + 10x^2 + 9}\dd x
	+ \int_{\gamma_R}\frac{x^2 - x + 2}{x^4 + 10x^2 + 9}\dd x
	= 2\pi i\left(\operatorname{Res}_{i}\left(\frac{x^2 - x + 2}
	{x^4 + 10x^2 + 9}\right)
	+ \operatorname{Res}_{3i}\left(\frac{x^2 - x + 2}
	{x^4 + 10x^2 + 9}\right)\right).
    \end{equation*}
    Noting that the roots of the denominator each have multiplicity $1$,
    i.e. are simple poles, the sum of the residues equals
    \begin{equation*}
	\frac{i^2 - i + 2}{(i+i)(i^2 + 9)}
	+ \frac{(3i)^2 - 3i + 2}{((3i)^2 + 1)(3i + 3i)}
	= \frac{1-i}{16i} + \frac{-7-3i}{-48i}
	= \frac 1{48i}(3 - 3i + 7 + 3i) = \frac 5{24i}.
    \end{equation*}
    Taking $R$ arbitrarily large,
    the latter integral on the right-hand side goes to zero because
    \begin{equation*}
	\int_{\gamma_R}\frac{x^2 - x + 2}{x^4 + 10x^2 + 9}\dd x
	= \int_0^{\pi}\frac{R^2e^{2i\theta} - Re^{i\theta} + 2}
	{R^4e^{4i\theta} + 10R^2e^{2i\theta} + 9}\cdot Re^{i\theta}\dd\theta
	= \int_0^\pi \frac 1{R}\cdot
	\frac{e^{2i\theta}-R^{-1}e^{i\theta}+1R^{-2}}
	{e^{4i\theta} + 10R^{-2}e^{2i\theta} + 9}\cdot e^{i\theta}\dd\theta,
    \end{equation*}
    which, in magnitude, is bounded above by
    \begin{equation*}
	\int_0^\pi \frac 1R\cdot\frac{1 + R^{-1} + 2R^{-2}}
	{1-10R^{-1}-9R^{-2}}\dd\theta
	= \frac \pi R\cdot O(1).
    \end{equation*}
    Thus, in the limit, we get that
    \begin{equation*}
	\int_{-\infty}^\infty \frac{x^2 - x + 2}{x^4 + 10x^2 + 9}\dd x
	= 2\pi i\cdot \frac 5{24i}
	= \boxed{\frac{5\pi}{12}}.
    \end{equation*}

    (d) Since the integrand is even,
    \begin{equation*}
	\int_0^\infty\frac{x^2}{(x^2 + a^2)^3}\dd x
	= \half \int_{-\infty}^\infty\frac{x^2}{(x^2 + a^2)^3}\dd x.
    \end{equation*}
    By the same procedure as with (c),
    \begin{equation*}
	\int_{-\infty}^\infty \frac{x^2}{(x^2 + a^2)^3}
	= 2\pi i\operatorname{Res}_{ai}\left(\frac{x^2}{(x^2+a^2)^3}\right).
    \end{equation*}
    Since $(x^2 + a^2)^3 = (x+ai)^3(x-ai)^3$,
    $ai$ is a pole of order $3$, so the residue is equal to
    \begin{equation*}
	\frac 1{2!}\cdot\frac{\dd^2}{\dd x^2}\frac{x^2(x-ai)^3}{(x^2+a^2)^3}
	\bigg|_{x=ai}
	= \half\left(\frac{12(ai)^2}{(2ai)^5}
	- \frac{12ai}{(2ai)^4} + \frac 2{(2ai)^3}\right)
	= \frac{-i}{16a^3},
    \end{equation*}
    meaning
    \begin{equation*}
	\int_{0}^\infty \frac{x^2}{(x^2 + a^2)^3}
	\half\int_{-\infty}^\infty \frac{x^2}{(x^2 + a^2)^3}
	= \half\cdot2\pi i\cdot \frac{-i}{16a^3}
	= \boxed{\frac{\pi}{16a^3}}.\qedhere
    \end{equation*}
\end{soln}

\pagebreak
\begin{problem}
    [p.136\#1]
    Show that $\abs{f(z)}\le 1$ for $\abs z\le 1$ implies
    \begin{equation*}
	\frac{\abs{f'(z)}}{1-\abs{f(z)}^2}
	\le \frac{1}{1-\abs{z}^2}.
    \end{equation*}
\end{problem}

\begin{soln}
    Using equation (36) on p.136 with
    $R = M = 1$, we get that for any $z_0$ in the unit disk,
    \begin{equation*}
	\abs{\frac{(f(z) - f(z_0))/(z-z_0)}{1 - \ol{f(z_0)}f(z)}}
	\le \abs{\frac{1}{1-\ol{z_0}z}}.
    \end{equation*}
    Taking the limit as $z\to z_0$ gives the desired inequality
    at $z_0$, since the denominators become
    $1 - \ol{f(z_0)}f(z_0)$ and $1-\ol z_0z_0$
    and the numerator on the left becomes $f'(z_0)$.
\end{soln}

\begin{problem}
    [p.136\#5]
    Show by use of Schwarz's lemma that
    every one-to-one conformal mapping of a disk onto another
    (or a half-plane) is given by a linear transformation.
\end{problem}

\begin{soln}
    Suppose first that $f$ sends the unit disk to itself
    biholomorphically. Then, if $a\in\Delta$ is the point
    such that $f(a) = 0$, compose $f$ with the linear transformation
    \begin{equation*}
        T\colon\Delta\to\Delta,\qquad
	Tz = \frac{a-z}{1-\ol az}
    \end{equation*}
    sending $0$ to $a$, so that $g = f\circ T$
    is an automorphism of $\Delta$ sending $0$ to $0$.
    Then let $h\colon \Delta\to\Delta$ be the inverse of $g$,
    and apply Schwarz's lemma to $g$ and $h$ to get that
    $\abs{g'(0)}$ and $\abs{h'(0)}$ are both at most $1$.
    By the chain rule $1 = g'(0)\cdot h'(g(0)) = g'(0)\cdot h'(0)$,
    implying that $\abs{g'(0)} = \abs{h'(0)} = 1$.
    The equality case of Schwarz's lemma tells us that
    $g(z) = cz$ for some constant $c$ with $\abs c = 1$,
    hence $f(z) = cT\inv(z) = cT(z)$, which is a linear transformation.

    In the general case, where $f\colon U\to V$
    is biholomorphic, the function
    $\varphi\circ f\circ\psi\inv\colon\Delta\to\Delta$
    where $\psi\colon U\to \Delta$ and $\psi\colon V\to\Delta$
    are linear transformations satisfies the hypotheses of the problem,
    so the prevoius paragraph tells us that it is a linear transformation,
    hence so is $f$.
\end{soln}

\begin{problem}
    [p.136\#6]
    If $\gamma$ is a piecewise differentiable arc contained in $\abs z < 1$
    the integral
    \begin{equation*}
	\int_\gamma \frac{\abs{\dd z}}{1-\abs z^2}
    \end{equation*}
    is called the \emph{noneuclidean} or \emph{hyperbolic} length of $\gamma$.
    Show that an analytic function $f(z)$ with $\abs{f(z)} < 1$
    for $\abs z < 1$ maps every $\gamma$ to
    an arc with smaller or equal noneuclidean length.
    Deduce that a linear transformation of a disk onto itself
    preserves noneuclidean lengths, and check the result
    by an explicit computation.
\end{problem}

\begin{soln}
    Parametrize $\gamma$ as a function $\gamma\colon[0,1]\to\Delta$.
    By p.136\#1, we have
    \begin{equation*}
	\frac{\abs{\gamma'(t)\cdot f'(\gamma(t))}}
	{1-\abs{f(\gamma(t))}^2}
	\le \frac{\abs{\gamma'(t)}}{1-\abs{\gamma(t)}^2}
    \end{equation*}
    for all $t\in[0,1]$, so integrating over $[0,1]$
    and doing a change of variables gives
    \begin{equation*}
	\int_{f(\gamma)}\frac{\abs{\dd z}}{1-\abs z^2}
	\le \int_{\gamma}\frac{\abs{\dd z}}{1-\abs z^2}.
    \end{equation*}
    If $T$ is a linear transformation on the unit disk then
    the hyperbolic length of $T\gamma$ is less than or equal
    to that of $\gamma$ and the hyperbolic length of $T\gamma$
    is less than or equal to that of $T\inv(T\gamma) = \gamma$,
    i.e. $T\gamma$ and $\gamma$ have the same hyperbolic length.

    Concretely, suppose that $f(z) = \frac{c(a-z)}{1-\ol az}$
    is an automorphism of $\Delta$, where $\abs c = 1$ and $\abs a < 1$.
    Then for a piecewise differentiable curve $\gamma\colon[0,1]\to\CC$,
    \begin{equation*}
	\frac{\abs{\gamma'(t)\cdot f'(\gamma(t))}}
	{1-\abs{f(\gamma(t))}^2}
	= \frac{\abs{\gamma'(t)}\cdot\abs{\frac{1-\abs a^2}
	{(1-\ol a\gamma(t))^2}}}{1 -
	\abs{\frac{a-\gamma(t)}{1-\ol a\gamma(t)}}^2}
	= \frac{\abs{\gamma'(t)}\cdot(1-\abs a^2)}
	{\abs{1-\ol a\gamma(t)}^2 - \abs{a-\gamma(t)}^2}.
    \end{equation*}
    The squared magnitudes in the denominator
    can be expanded by writing them as the argument
    times its conjugate to get
    \begin{equation*}
	(1-\ol a\gamma(t))(1- a\ol{\gamma(t)})
	- (a-\gamma(t))(\ol a-\ol{\gamma(t)})
	= 1 - \abs a^2 - \abs{\gamma(t)}^2
	+ \abs a^2\abs{\gamma(t)}^2
	= (1-\abs a^2)(1-\abs{\gamma(t)}^2)
    \end{equation*}
    so that the fraction above equals
    \begin{equation*}
	\frac{\abs{\gamma'(t)\cdot f'(\gamma(t))}}
	{1-\abs{f(\gamma(t))}^2}
	= \frac{\abs{\gamma'(t)}}{1-\abs{\gamma(t)}^2}.
    \end{equation*}
    Thus,
    \begin{equation*}
	\int_\gamma\frac{\abs{dz}}{1-\abs z^2}
	= \int_0^1 \frac{\abs{\gamma'(t)}}{1-\abs{\gamma(t)}^2}\dd t
	= \int_0^1 \frac{\abs{\gamma'(t)\cdot f'(\gamma(t))}}
	{1-\abs{f(\gamma(t))}^2}\dd t
	= \int_{f(\gamma)}\frac{\abs{\dd z}}{1 - \abs z^2}
    \end{equation*}
    as desired.
\end{soln}

\begin{problem}
    [p.186\#1]
    Prove that the Laurent development is unique.
\end{problem}

\begin{soln}
    This follows from the fact that the coefficients
    $a_n$ in the Laurent series
    \begin{equation*}
	f(z) = \sum_{n=-\infty}^\infty a_n(z-z_0)^n
    \end{equation*}
    can be directly computed in terms of $f$, $z_0$, and $n$ as
    \begin{equation*}
	a_n = \frac 1{2\pi i}\int_{\abs{z-z_0}=r}\frac{f(\zeta)}
	{(\zeta-z_0)^{n+1}}\dd\zeta,
    \end{equation*}
    where $r$ is between the smaller and larger radius
    of the annulus of convergence of the series.
    Note also that the integral is independent of
    the particular choice of $r$ by homotopy invariance/Green's theorem.
\end{soln}










\end{document}
